\documentclass[a4paper,12pt,liststotocnumbered]{scrartcl}

\usepackage[utf8x]{inputenc}
\usepackage{ngerman}
\usepackage{tabularx}
\usepackage{graphicx}
\usepackage{hyperref}
\usepackage{calc}
\usepackage{scrpage2}

\pagestyle{scrheadings}

\hypersetup{colorlinks=false,pdfborder=0 0 0}

% command for version
\newcommand{\lxVersion}[0]{0.6}
% command for date
\newcommand{\lxDate}[0]{\today}

\chead[]{}
\lohead[Webtechnologie I - T09]{Webtechnologie I - T09}
\rohead[\today]{\today}

%\chead[{
%\lohead[Gateway- / Proxyserver]{Gateway- / Proxyserver}
%\rohead[Projektname: ProximusPrime]{Projektname: ProximusPrime}
%}]

\title{Anforderungsdokument Fotoalbum}
\author{Sebastian Eberlein, Sven Salzwedel - Team 09}

\begin{document}

\begin{figure}[h]
	\begin{center}
		\includegraphics[width=\textwidth/2]{logo}
	\end{center}
\end{figure}

\begin{tabularx}{\textwidth}{lX}
	Fach:&Webtechnologie I\\
	Semester:&Sommersemester 2007\\
	Prüfungsleistung:&Praktischer Leistungsnachweis\\
	&\\
	Teammitglieder:&Sebastian Eberlein\\
	&Sven Salzwedel\\
	Teamnummer:&T09\\
\end{tabularx}

\begin{center}
\end{center}

\begin{center}
	\Large{\textbf{Thema:} Webbasiertes Fotoalbum}\\
\end{center}
\begin{center}
	\Large{\textbf{Notiz zum Anforderungsdokument}}\\
\end{center}

\begin{center}
	Datum: \today\\
\end{center}

\newpage

%\maketitle

\section{Funktionsänderungen}

\subsection{Albumschlüssel}

Es wurde eine weitere Funktion zum Programm hinzugefügt, welche nicht im
ursprünglichen Anforderungsdokument erwähnt wurde. Es kann für ein Album
ein Schlüssel vom Administrator generiert werden, sodass das jeweilige
Album nur noch von den Personen angesehen werden kann, welche den
dazugehörigen Schlüssel wissen.

\subsection{Thumbnailgenerierung}

Eine weitere Abhängigkeit die nicht im Anforderungsdokument beschrieben
wurde ist eine zusätzliche Bibliothek \emph{JIU}
(\url{http://schmidt.devlib.org/jiu/index.html}), die eingebunden werden
musste um die Erzeugung der Thumbnails zu ermöglichen.

\subsection{XSLT Prozessor}

Der im Anforderungsdokument beschriebene XSLT Prozessor wurde nicht
implementiert.  Dessen Aufgabe übernehmen die Clients selbst. Die
xml-Dateien werden direkt an den Browser ausgeliefert, welcher dann die
xslt-Dateien nachlädt und so die Seiten selbst generiert. Da die in den
Vorgaben geforderte Clients diesen Mechanismus unterstützen, bleiben die
geplanten Funktionalitäten genauso vorhanden als würde man einen
serverseitigen XSLT-Prozessor einsetzen.

\newpage

\end{document}
